\documentclass[../../main.tex]{subfiles}%
\graphicspath{{\subfix{./image/}}{\subfix{../../image/}}}

\begin{document}

\section{参考文献与注记}\label{section:ch05-12}

关于 $\BV$ 理论，我们主要依据 Giusti\cite{giusti1984} 与 Federer\cite{federer1969} 第 4.5 节，并参考 Simon\cite{simon1984} 第 6 节。结构定理（\cref{theorem:ch05-5-1}）例如在\cite{simon1984} 第 6.1 节中给出陈述。% XREF p=255 "The Structure Theorem is stated, for instance, in [S, Section 6.1]"
\cref{section:ch05-02}中的下半连续性定理取自\cite{giusti1984} 第 1.9 节，局部逼近定理取自\cite{giusti1984} 中编号为 1.17 的定理（此结果应归于 Anzellotti 与 Giaquinta）。% XREF p=255 "The Lower Semicontinuity Theorem in Section 5.2 is [G, Section 1.9], and the Local Approximation Theorem is [G, Theorem 1.17]. (This result is due to Anzellotti and Giaquinta)"
\cref{section:ch05-02}中的紧性断言取自\cite{giusti1984} 中编号为 1.19 的定理。% XREF p=255 "The compactness assertion in Section 5.2 follows [G, Theorem 1.19]"
\cref{section:ch05-03}中关于迹的讨论遵循\cite{giusti1984} 第 2 章。% XREF p=255 "The discussion of traces in Section 5.3 follows [G, Chapter 2]"
\cref{section:ch05-04}中对延拓的处理是\cite{giusti1984} 中编号为 2.13 的注记的详细展开。% XREF p=255 "Our treatment of extensions in Section 5.4 is an elaboration of [G, Remark 2.13]"

$\BV$ 函数的余面积公式（\cref{theorem:ch05-5-9}）应归于 Fleming 与 Rishel\cite{fleming1960}，其证明按照\cite{giusti1984} 中编号为 1.23 的定理进行。% XREF p=255 "The coarea formula for BV functions, due to Fleming and Rishel [Fl-R], is proved as in [G, Theorem 1.23]"
关于等周不等式，参看\cite{giusti1984} 中编号为 1.28 的定理与编号为 1.29 的推论。\cref{section:ch05-06}中的注记与\cite{federer1969} 第 4.5.9 节之 (18) 有关。% XREF p=255 "The remark in Section 5.6 is related to [F, Section 4.5.9(18)]"
\cref{theorem:ch05-5-12}应归于 Fleming；我们遵循\cite{federer1972}。% XREF p=255 "Theorem 5.12 is due to Fleming; we followed [F-Z]"
\cref{section:ch05-07}与\cref{section:ch05-08}中关于约化边界与测度论边界的结果来自\cite{giusti1984} 第 3、4 章；这些断言最初由 De Giorgi 建立。% XREF p=255 "The results in Sections 5.7 and 5.8 on the reduced and measure-theoretic boundaries are from [G, Chapters 3 and 4]" | p=255 "The results in Sections 5.7 and 5.8 on the reduced and measure-theoretic boundaries are from [G, Chapters 3 and 4]"

Federer\cite{federer1969} 第 4.5.9 节列出了 $\BV$ 函数的一长串性质，我们从中提取了\cref{section:ch05-09}中所述的这套理论。% XREF p=255 "Federer [F, Section 4.5.9] presents a long list of properties of BV functions, from which we extracted the theory set forth in Section 5.9"
本质变差出现在\cite{federer1969} 第 4.5.10 节中，而\cref{section:ch05-11}中所述的有限周长判别准则即\cite{federer1969} 第 4.5.11 节。% XREF p=255 "Essential variation occurs in [F, Section 4.5.10] and the criterion for finite perimeter described in Section 5.11 is [F, Section 4.5.11]"

L. Ambrosio 与 E. De Giorgi\cite{ambrosio1988} 引进了「特殊」有界变差函数类，记作 $\mathrm{SBV}$，其梯度的奇异部分支撑在跳跃集 $J$ 上。另见 Ambrosio\cite{ambrosio1995}。

\end{document}
